\section*{الفصل \ref{ch:g12:mechanical-waves} --- الأمواج الميكانيكية}

\begin{solution}{exo:g12:mechanical-waves:1}
(أ) $v = 320/40 = \qty{8.0}{m/s}$. (ب) لا يتحرك أي متفرج على امتداد
الحلقة: فالمنتقل هو الاضطراب (أي فعل الوقوف)، لا الجمهور.
(ج) $12/1.5 = \qty{8.0}{m/s}$ --- وهو متسق.
\end{solution}

\begin{solution}{exo:g12:mechanical-waves:2}
(أ) \omterm{def:g12:mechanical-waves:transverse}{مستعرضة}؛ (ب) \omterm{def:g12:mechanical-waves:transverse}{طولية}؛ (ج) \omterm{def:g12:mechanical-waves:transverse}{طولية}؛ (د) \omterm{def:g12:mechanical-waves:transverse}{مستعرضة}؛
(ه) P \omterm{def:g12:mechanical-waves:transverse}{طولية}، وS \omterm{def:g12:mechanical-waves:transverse}{مستعرضة}.
\end{solution}

\begin{solution}{exo:g12:mechanical-waves:3}
(أ) $d = 340 \times 4.0 \approx \qty{1.4}{km}$.
(ب) $1500 \times 4.0 = \qty{6.0}{km}$. (ج) \omterm{def:g12:mechanical-waves:speed}{سرعة الانتشار} في
الوسط المقطوع.
\end{solution}

\begin{solution}{exo:g12:mechanical-waves:4}
$\lambda = v/f$: أي $340/440 = \qty{0.77}{m}$ في الهواء،
و$1500/440 = \qty{3.4}{m}$ في الماء. ويبقى \omterm{def:g10:signals-and-waves:frequency}{التواتر} ثابتًا (فالمنبع
يمليه)؛ أما $v$ و$\lambda$ فيتغيران كلاهما.
\end{solution}

\begin{solution}{exo:g12:mechanical-waves:5}
$\lambda = \qty{12}{cm}$؛ $T = \lambda/v = 0.12/0.30 = \qty{0.40}{s}$؛
$f = 1/T = \qty{2.5}{Hz}$.
\end{solution}

\begin{solution}{exo:g12:mechanical-waves:6}
(أ) $v = 2.4/0.30 = \qty{8.0}{m/s}$. (ب) ومن الصورة الأولى يبقى على
القمة $\qty{6.0}{m}$: أي $t = 6.0/8.0 = \qty{0.75}{s}$ بعدها.
(ج) ويرتفع الشريط ثم ينخفض، عرضيًّا خالصًا، ويبقى
عند $x = \qty{5.0}{m}$: فالحبل لا ينتقل.
\end{solution}

\begin{solution}{exo:g12:mechanical-waves:7}
(أ) سرعة الاقتراب \qty{4.0}{m/s}، والفجوة \qty{4.0}{m}: ومنه $t = \qty{1.0}{s}$.
(ب) بالتراكب: $3.0 + 2.0 = \qty{5.0}{cm}$.
(ج) $3.0 - 2.0 = \qty{1.0}{cm}$. (د) تخرج نبضتان سليمتان وتواصلان
سيرهما دون تغير، كأنهما لم تلتقيا قط.
\end{solution}

\begin{solution}{exo:g12:mechanical-waves:8}
(أ) الرفيع: فكتلة أقلّ في \omterm{def:g10:orders-of-magnitude:unit}{المتر} عند \omterm{def:g10:inertia:inventory}{الشدّ} نفسه تعني \omterm{def:g10:signals-and-waves:wave}{موجة}
أسرع. (ب) تزداد (لأن \omterm{def:g10:inertia:inventory}{الشدّ} أعلى). (ج) لا: فالسرعة لا تتوقف
إلا على الوسط، لا على حجم النبضة.
\end{solution}

\begin{solution}{exo:g12:mechanical-waves:9}
السكة: $1800/5900 \approx \qty{0.31}{s}$؛ والهواء:
$1800/340 \approx \qty{5.3}{s}$؛ والفارق $\approx \qty{5.0}{s}$.
\end{solution}

\begin{solution}{exo:g12:mechanical-waves:10}
بأخذ $\lambda = \qty{24}{cm}$: (أ) $2\lambda$، \omterm{def:g12:mechanical-waves:phase}{فعلى توافق الطور}؛
(ب) $\lambda/2$، فعلى \omterm{def:g12:mechanical-waves:phase}{تعاكس الطور}؛ (ج) $\tfrac32\lambda$، فعلى \omterm{def:g12:mechanical-waves:phase}{تعاكس الطور}؛
(د) $1.25\lambda$، فلا هذا ولا ذاك؛ (ه) $\lambda/4$، فلا هذا ولا ذاك.
\end{solution}

\begin{solution}{exo:g12:mechanical-waves:11}
$\lambda_P = 6.0 \times 2.0 = \qty{12}{km}$؛
$\lambda_S = 3.5 \times 2.0 = \qty{7.0}{km}$ --- أي مئات أضعاف
حجم بناية. والنقاط الأقرب بكثير من $\lambda$ تكون شبه متوافقة
الطور، فيتحرك الحيّ كله معًا.
\end{solution}

\begin{solution}{exo:g12:mechanical-waves:12}
(أ) $T = \qty{0.50}{s}$، $f = \qty{2.0}{Hz}$.
(ب) $\lambda = \qty{0.75}{m}$.
(ج) $v = \lambda/T = 0.75/0.50 = \qty{1.5}{m/s}$.
(د) عند $t = T = \qty{0.50}{s}$. (ه) $1.875/0.75 = 2.5$ طول \omterm{def:g10:signals-and-waves:wave}{موجة}
$= (2 + \tfrac12)\lambda$: أي \omterm{def:g12:mechanical-waves:phase}{تعاكس الطور}.
\end{solution}

\begin{solution}{exo:g12:mechanical-waves:13}
(أ) $T = 80/10 = \qty{8.0}{s}$، $f = \qty{0.125}{Hz}$.
(ب) $v = \lambda/T = 120/8.0 = \qty{15}{m/s} = \qty{54}{km/h}$.
(ج) $t = 900/15 = \qty{60}{s}$. (د) لا: فالأمواج لا تنقل مادة ---
والعوّامة الراسية لا تفعل إلا أن ترتفع وتنخفض.
\end{solution}

\begin{solution}{exo:g12:mechanical-waves:14}
(أ) بالاقتراب بسرعة \qty{6.0}{m/s} على مدى \qty{6.0}{m}: $t = \qty{1.0}{s}$؛
ويتلاشى الانتقالان المتعاكسان --- فيصير الحبل مسطّحًا للحظة.
(ب) لا: فنقاطه متحركة (إذ تتجمع السرعات العرضية ولا
تتلاشى)؛ \omterm{def:g10:energy-conservation:energy}{والطاقة} كلها حركية. (ج) وتعاود النبضتان الظهور
وتواصلان سيرهما دون تغير. (د) فالحبل المسطّح يحمل سرعة، ومن ثم \omterm{def:g10:energy-conservation:energy}{طاقة}:
\omterm{def:g10:signals-and-waves:wave}{فالموجة} مخزونة للحظة في الحركة، لا في الشكل.
\end{solution}

\begin{solution}{exo:g12:mechanical-waves:15}
(أ) $2\pi r$: أي \qty{3.1}{m} و\qty{50}{m} --- ومنه $\times 16$.
(ب) وتهبط \omterm{def:g10:energy-conservation:energy}{الطاقة} في \omterm{def:g10:orders-of-magnitude:unit}{المتر} بالعامل $16$، فتهبط \omterm{def:g10:signals-and-waves:amplitude}{السعة} (لأن
\omterm{def:g10:energy-conservation:energy}{الطاقة} تنمو مع \omterm{def:g10:signals-and-waves:amplitude}{السعة}). (ج) والجبهات المستقيمة لا
تطول: فلا يبقى إلا الامتصاص. (د) والوسط نشط: فكل
متفرج يقف على \omterm{def:g10:energy-conservation:energy}{طاقة} عضلاته هو، فيعيد تزويد \omterm{def:g10:signals-and-waves:wave}{الموجة}.
\end{solution}

\begin{solution}{pb:g12:mechanical-waves:1}
\textbf{1.} تضغط أمواج P في جهة الانتقال: فهي \omterm{def:g12:mechanical-waves:transverse}{طولية}؛ وتهزّ أمواج S
جانبيًّا: فهي \omterm{def:g12:mechanical-waves:transverse}{مستعرضة}.

\textbf{2.} تصل P، وهي الأسرع، أولًا في كل مكان.
$t_P = 120/6.0 = \qty{20}{s}$؛ $t_S = 120/3.5 \approx \qty{34}{s}$.

\textbf{3.} $\Delta t = d/v_S - d/v_P = d\,(1/v_S - 1/v_P)$؛ ومن أجل
\qty{120}{km}: $34.3 - 20 \approx \qty{14}{s}$.

\textbf{4.} حُلّ لإيجاد $d$: $d = \dfrac{v_P v_S}{v_P - v_S}\,\Delta t
= \dfrac{6.0 \times 3.5}{2.5}\,\Delta t = \qty{8.4}{km}$ لكل \omterm{def:g10:orders-of-magnitude:unit}{ثانية}
من الفارق.

\textbf{5.} $\Delta t \propto d$: فتسجيل محطة واحدة يعطي
مسافتها إلى مركز الزلزال دون معرفة زمن النشوء.

\textbf{6.} $d_A = 8.4 \times 25 = \qty{2.1e2}{km}$.

\textbf{7.} يعطي الفارق مسافة لا اتجاهًا: فيقع مركز الزلزال
في مكان ما على الدائرة ذات نصف القطر $d_A$ حول A.

\textbf{8.} $d_B = 8.4 \times 15 = \qty{1.3e2}{km}$؛
$d_C = 8.4 \times 32 \approx \qty{2.7e2}{km}$.

\textbf{9.} تلتقي دائرتان في نقطتين؛ وتنتقي الثالثة إحداهما:
فالتقاطع المشترك هو مركز الزلزال.

\textbf{10.} $t_P = 210/6.0 = \qty{35}{s}$: فقد بدأ الزلزال عند
$09{:}14{:}32 - \qty{35}{s} = 09{:}13{:}57$.

\textbf{11.} $\qty{200}{m/s} = \qty{720}{km/h}$ --- أي طائرة ركاب نفّاثة
في إبحارها.

\textbf{12.} $t = \num{6.0e5}/200 = \qty{3.0e3}{s} = \qty{50}{min}$.

\textbf{13.} $t_P = 600/6.0 = \qty{100}{s}$. والإنذار:
$3000 - 100 = \qty{2.9e3}{s} \approx \qty{48}{min}$.

\textbf{14.} $\lambda = vT = 200 \times 900 = \qty{1.8e5}{m} =
\qty{180}{km}$: أي ارتفاع نحو \omterm{def:g10:orders-of-magnitude:unit}{متر} ممتدًّا على \qty{180}{km} ---
فتتسلق السفينة ميلًا غير محسوس على مدى دقائق عدة.

\textbf{15.} $\lambda = 10 \times 900 = \qty{9.0}{km}$: فتتكدّس \omterm{def:g10:signals-and-waves:wave}{الموجة}
عشرين ضعفًا، وتنحشر طاقتها في طول أقصر، فيجب أن ينمو
ارتفاعها --- فينتصب التسونامي عند الشاطئ.

\textbf{16.} $2\pi r$: أي \qty{63}{km}، ثم $\approx \qty{1.3e3}{km}$
--- ومنه $\times 21$.

\textbf{17.} تُتقاسم \omterm{def:g10:energy-conservation:energy}{الطاقة} نفسها على جبهة أطول $21\times$
مرة، فتهبط \omterm{def:g10:signals-and-waves:amplitude}{السعة} حتى في قشرة بلا ضياع؛ وتضيف القشرة \omterm{def:g11:lenses-and-eye:images}{الحقيقية}
امتصاصًا.

\textbf{18.} أمواج S، وهي قصّ مستعرض، لا تقدر أن تعبر سائلًا: فظلّ
أمواج S يبيّن أن \omterm{def:g11:fundamental-interactions:ladder}{النواة} (الخارجية) سائلة.

\textbf{19.} $\num{1e15}/\num{1.5e9} \approx \num{6.7e5}$: أي نحو
سبعمئة ألف قطار بكامل سرعتها.

\textbf{20.} وقع الزلزال عند $09{:}13{:}57$، في عرض البحر عند
تقاطع الدوائر الثلاث، على بعد \qty{210}{km} من المحطة A. وانطلق
إنذار \omterm{def:g10:signals-and-waves:wave}{موجة} P في الميناء بعد \qty{100}{s} من النشوء، أي نحو
\qty{48}{min} قبل وصول التسونامي.
\end{solution}
